\clearpage
\SourceLeaf{613}
\ArithmeticHead{584}{L}{INHALT.}
\begingroup
\newcommand{\ContentsRow}[5]{\par\noindent\hangindent=#1\hspace*{#1}#2\hspace{.4em}\mbox{\footnotesize #3}\hfill\mbox{#4\quad #5}\par}
\ContentsRow{0pt}{7.\quad Kettenbrüche, 1796 Mai 24}{.\,.\,.\,.\,.\,.\,.\,.\,.\,.\,.\,.\,.\,.\,.\,.\,.\,.\,.\,.\,.\,.\,.}{Seite}{490}
\ContentsRow{0pt}{8.\quad Rekurrente Reihen, 1796 Mai 26}{.\,.\,.\,.\,.\,.\,.\,.\,.\,.\,.\,.\,.\,.\,.\,.\,.\,.\,.\,.\,.}{—}{493}
\ContentsRow{0pt}{9.\quad Asymptotische Gesetze der Zahlentheorie, 1796 Mai 31}{.\,.\,.\,.\,.\,.\,.\,.\,.\,.\,.\,.\,.}{—}{493}
\ContentsRow{0pt}{10.\quad Rekurrente Reihen, 1796 Juni 3}{.\,.\,.\,.\,.\,.\,.\,.\,.\,.\,.\,.\,.\,.\,.\,.\,.\,.\,.\,.\,.\,.}{—}{494}
\ContentsRow{0pt}{11, 12.\quad Teiler und Perioden von Zahlen, 1796 Juni 5}{.\,.\,.\,.\,.\,.\,.\,.\,.\,.\,.\,.\,.\,.\,.\,.}{—}{494}
\ContentsRow{0pt}{13, 14.\quad Asymptotische Gesetze der Zahlentheorie, 1796 Juni 19, 20}{.\,.\,.\,.\,.\,.\,.\,.\,.\,.\,.}{—}{495}
\ContentsRow{0pt}{15.\quad Binäre quadratische Formen, 1796 Juni 22}{.\,.\,.\,.\,.\,.\,.\,.\,.\,.\,.\,.\,.\,.\,.\,.\,.\,.}{—}{496}
\ContentsRow{0pt}{16.\quad Zweiter Beweis des Reziprozitätssatzes der quadratischen Reste, 1796 Juni 27}{.\,.\,.\,.\,.}{—}{496}
\ContentsRow{0pt}{17, 18.\quad Ternäre quadratische Formen, 1796 Juli 3, 10}{.\,.\,.\,.\,.\,.\,.\,.\,.\,.\,.\,.\,.\,.\,.}{—}{496,497}
\ContentsRow{0pt}{19.\quad Binäre quadratische Formen, 1796 Juli}{.\,.\,.\,.\,.\,.\,.\,.\,.\,.\,.\,.\,.\,.\,.\,.\,.\,.\,.}{—}{497}
\ContentsRow{0pt}{20.\quad Rekurrente Reihen, 1796 Juli 16}{.\,.\,.\,.\,.\,.\,.\,.\,.\,.\,.\,.\,.\,.\,.\,.\,.\,.\,.\,.\,.}{—}{498}
\ContentsRow{0pt}{21.\quad Algebraische Kurven, 1796 Juli 31}{.\,.\,.\,.\,.\,.\,.\,.\,.\,.\,.\,.\,.\,.\,.\,.\,.\,.\,.\,.\,.}{—}{498}
\ContentsRow{0pt}{22.\quad Binomische Kongruenzen, 1796 August 3}{.\,.\,.\,.\,.\,.\,.\,.\,.\,.\,.\,.\,.\,.\,.\,.\,.\,.}{—}{498}
\ContentsRow{0pt}{23.\quad Zum 3.\ und 4.\ Beweise des Reziprozitätssatzes der quadratischen Reste, 1796 August 13}{.}{—}{499}
\ContentsRow{0pt}{24.\quad Komplexe Größen, 1796 August 14}{.\,.\,.\,.\,.\,.\,.\,.\,.\,.\,.\,.\,.\,.\,.\,.\,.\,.\,.\,.}{—}{499}
\ContentsRow{0pt}{25.\quad Zahlentheoretisches (vielleicht zu Nr.\ 23 gehörig), 1796 August 16}{.\,.\,.\,.\,.\,.\,.\,.\,.}{—}{499}
\ContentsRow{0pt}{26.\quad Kongruenzen, 1796 August 18}{.\,.\,.\,.\,.\,.\,.\,.\,.\,.\,.\,.\,.\,.\,.\,.\,.\,.\,.\,.\,.\,.}{—}{499}
\ContentsRow{0pt}{27.\quad Unbestimmte Gleichung ersten Grades, 1796 August 19}{.\,.\,.\,.\,.\,.\,.\,\hspace{.55em}.\,.\,.\,.\,.}{—}{500}
\ContentsRow{0pt}{28.\quad Potenzsummen der Wurzeln, 1796 August 21}{.\,.\,.\,.\,.\,.\,.\,.\,.\,.\,.\,.\,.\,.\,.\,.\,.}{—}{500}
\ContentsRow{0pt}{29.\quad Sinus bezw.\ Kosinus höherer Ordnung, 1796 August 21}{.\,.\,.\,.\,.\,.\,.\,.\,.\,.\,.\,.\,.}{—}{501}
\ContentsRow{0pt}{30.\quad Zum 3.\ und 4.\ Beweise des Reziprozitätssatzes der quadratischen Reste, 1796 September 2}{.}{—}{501}
\ContentsRow{0pt}{31.\quad Asymptotisches Gesetz der Zahlentheorie, 1796 September 6}{.\,.\,.\,.\,.\,.\,.\,.\,.\,.\,.\,.}{—}{502}
\ContentsRow{0pt}{32, 33.\quad Umkehrung von Integralen, 1796 September 9 und 14}{.\,.\,.\,.\,.\,.\,.\,.\,.\,.\,.\,.}{—}{502}
\ContentsRow{0pt}{34.\quad \textsc{Tschirnhausen}sche Transformation, 1796 September 16}{.\,.\,.\,.\,.\,.\,.\,.\,.\,.\,.\,.\,.}{—}{503}
\ContentsRow{0pt}{35.\quad Umformung von Funktionen der Gleichungswurzeln. 1796 September 16}{.\,.\,.\,.\,.\,.\,.}{—}{503}
\ContentsRow{0pt}{36.\quad Elimination, 1796 September 16}{.\,.\,.\,.\,.\,.\,.\,.\,.\,.\,.\,.\,.\,.\,.\,.\,.\,.\,.\,.\,.\,.}{—}{503}
\ContentsRow{0pt}{37.\quad \textsc{Lagrange}sche Resolvente, 1796 September 17}{.\,.\,.\,.\,.\,.\,.\,.\,.\,.\,.\,.\,.\,.\,.\,.}{—}{503}
\ContentsRow{0pt}{38.\quad Einheitswurzeln, 1796 September 29}{.\,.\,.\,.\,.\,.\,.\,.\,.\,.\,.\,.\,.\,.\,.\,.\,.\,.\,.\,.}{—}{505}
\ContentsRow{0pt}{39.\quad Kubische Resolvente der Kreisteilungsgleichung, 1796 Oktober 1}{.\,.\,.\,.\,.\,.\,.\,.\,.\,.}{—}{505}
\ContentsRow{0pt}{40.\quad Irreduzibilität der Kreisteilungsgleichung für Primzahlgrad, 1796 Oktober 9}{.\,.\,.\,.\,.\,.}{—}{507}
\ContentsRow{0pt}{41.—43.\quad Über algebraische Gleichungen, 1796 Oktober 16, 18, 21}{.\,.\,.\,.\,.\,.\,.\,.\,.\,.\,.}{—}{507}
\ContentsRow{0pt}{44.\quad Interpolationsformel, 1796 November 25}{.\,.\,.\,.\,.\,.\,.\,.\,.\,.\,.\,.\,.\,.\,.\,.\,.\,.\,.}{—}{508}
\ContentsRow{0pt}{45.\quad Eine Potenzreihe, 1796 November 26}{.\,.\,.\,.\,.\,.\,.\,.\,.\,.\,.\,.\,.\,.\,.\,.\,.\,.\,.\,.}{—}{508}
\ContentsRow{0pt}{46.\quad Trigonometrische Formeln, 1796 Dezember}{.\,.\,.\,.\,.\,.\,.\,.\,.\,.\,.\,.\,.\,.\,.\,.\,.\,.}{—}{508}
\ContentsRow{0pt}{47.\quad Allgemeine Differentiation, 1796 Dezember 23}{.\,.\,.\,.\,.\,.\,.\,.\,.\,.\,.\,.\,.\,.\,.\,.\,.}{—}{508}
\ContentsRow{0pt}{48.\quad Quadratur parabolischer Kurven, 1796 Dezember 26}{.\,.\,.\,.\,.\,.\,.\,.\,.\,.\,.\,.\,.\,.}{—}{509}
\ContentsRow{0pt}{49.\quad \textsc{Lagranges} Umkehrungssatz bewiesen, 1796 Dezember 27}{.\,.\,.\,.\,.\,.\,.\,.\,.\,.\,.\,.}{—}{509}
\ContentsRow{0pt}{50.\quad Integralformeln, 1797 Januar 7}{.\,.\,.\,.\,.\,.\,.\,.\,.\,.\,.\,.\,.\,.\,.\,.\,.\,.\,.\,.\,.\,.}{—}{509}
\ContentsRow{0pt}{51.\quad Lemniskatische Funktionen begonnen, 1797 Januar 8}{.\,.\,.\,.\,.\,.\,.\,.\,.\,.\,.\,.\,.\,.}{—}{510}
\ContentsRow{0pt}{52.—54.\quad Integralrechnung, 1797 Januar 10, 12}{.\,.\,.\,.\,.\,.\,.\,.\,.\,.\,.\,.\,.\,.\,.\,.\,.\,.}{—}{510,511}
\ContentsRow{0pt}{55.\quad Konstruktion der regelmäßigen Vielecke, 1797 Januar 19}{.\,.\,.\,.\,.\,.\,.\,.\,.\,.\,.\,.\,.}{—}{511}
\ContentsRow{0pt}{56.\quad Restcharaktere von $-1$, $\pm2$, 1797 Februar 4}{.\,.\,.\,.\,.\,.\,.\,.\,.\,.\,.\,.\,.\,.\,.\,.\,.}{—}{512}
\ContentsRow{0pt}{57.\quad Quadratische ternäre Formen, 1797 Februar 6}{.\,.\,.\,.\,.\,.\,.\,.\,.\,.\,.\,.\,.\,.\,.\,.\,.}{—}{512}
\ContentsRow{0pt}{58.\quad Kettenbruch; Potenzreihe, deren Exponenten eine arithmetische Reihe zweiter Ordnung bilden, 1797 Februar 16}{.\,.\,.\,.\,.\,.\,.\,.\,.\,.\,.\,.\,.\,.\,.\,.\,.\,.\,.\,.\,.\,.\,.\,.\,.\,.}{—}{513}
\endgroup
